\section*{Hoofdstuk \ref{ch:g12:internal-heat} --- De inwendige warmte van de Aarde}

\begin{solution}{exo:g12:internal-heat:1}
(a) de \omterm{def:g12:internal-heat:heat-flow}{warmtestroom}; (b) de \omterm{def:g12:internal-heat:heat-production}{radiogene warmteproductie}; (c) \omterm{def:g12:internal-heat:secular-cooling}{seculaire afkoeling};
(d) \omterm{def:g12:internal-heat:convection}{thermische convectie}; (e) \omterm{def:g12:internal-heat:geothermal}{geothermische energie}.
\end{solution}

\begin{solution}{exo:g12:internal-heat:2}
(1) \qty{0.0916}{W/m^2}. (2) \qty{4.67e13}{W}. (3) \qty{0.029}{K/m}: een
verschil van één graad Celsius is een verschil van één kelvin.
\end{solution}

\begin{solution}{exo:g12:internal-heat:3}
Thorium-232, uranium-238, kalium-40 en uranium-235. KamLAND ziet alleen de
geoneutrino's van thorium-232 en uranium-238.
\end{solution}

\begin{solution}{exo:g12:internal-heat:4}
$490/\sqrt{1} = \qty{490}{mW/m^2}$; $490/2 = \qty{245}{mW/m^2}$;
$490/5 = \qty{98}{mW/m^2}$; $490/7 = \qty{70}{mW/m^2}$, ongeveer het gemiddelde van de
continenten.
\end{solution}

\begin{solution}{exo:g12:internal-heat:5}
Dicht bij het oppervlak volgt de temperatuur van de grond de dagen en de
seizoenen; pas onder een paar meter blijft ze op de gemiddelde temperatuur van het
jaar, en pas nog dieper toont ze de warmte die van beneden komt.
Boren woelt het gesteente en het water in het gat om en verandert hun
temperatuur, dus wacht je tot het gat is teruggekeerd naar de temperatuur van
het gesteente.
\end{solution}

\begin{solution}{exo:g12:internal-heat:6}
(1) $(265 - 10)/9.1 \approx \qty{28.0}{\celsius}$ per kilometer, dat is
\qty{0.0280}{K/m}. (2) $q = 2.72 \times 0.0280 \approx \qty{0.0762}{W/m^2}$,
ongeveer \qty{76}{mW/m^2}. (3) Iets boven het gemiddelde van de continenten,
\qty{70.9}{mW/m^2}, met ongeveer \qty{7}{\percent}: een gewone continentale waarde.
\end{solution}

\begin{solution}{exo:g12:internal-heat:7}
Continenten: $0.0709 \times \num{2.073e14} \approx \qty{1.47e13}{W} =
\qty{14.7}{TW}$. Oceanen: $0.1054 \times \num{3.028e14} \approx
\qty{3.19e13}{W} = \qty{31.9}{TW}$. Totaal ongeveer \qty{46.6}{TW}, waarvan
$31.9/46.6 \approx \qty{68}{\percent}$, twee derde, door de oceaanbodem
weggaat.
\end{solution}

\begin{solution}{exo:g12:internal-heat:8}
(1) $\num{1.88e-6} \times 8000 \approx \qty{0.0150}{W/m^2} = \qty{15.0}{mW/m^2}$.
(2) $15.0/70.9 \approx \qty{21}{\percent}$, een vijfde van de \omterm{def:g12:internal-heat:heat-flow}{warmtestroom}.
\end{solution}

\begin{solution}{exo:g12:internal-heat:9}
$\sqrt{t} = 490/70.9 \approx 6.91$, dus $t \approx 47.8$ miljoen jaar: oceaanbodem
die ouder is dan ongeveer 48 miljoen jaar, verliest minder warmte dan het gemiddelde
continent.
\end{solution}

\begin{solution}{exo:g12:internal-heat:10}
Uranium-235: $2^{1.4/0.704} = 2^{1.99} \approx 4$ keer meer. Kalium-40:
$2^{2.5/1.248} = 2^{2.00} \approx 4$ keer meer. Beide liggen twee halveringstijden terug.
\end{solution}

\begin{solution}{exo:g12:internal-heat:11}
(1) $10^2/10^{-6} = \qty{1e8}{s}$, ongeveer 3 jaar. (2) $(10^5)^2/10^{-6} =
\qty{1e16}{s}$, ongeveer 320 miljoen jaar. (3) $(\num{2.891e6})^2/10^{-6}
\approx \qty{8.4e18}{s}$, ongeveer 265 miljard jaar. De tijd groeit met het
kwadraat van de afstand: tien keer dikker, honderd keer langer.
\end{solution}

\begin{solution}{exo:g12:internal-heat:12}
Vermenigvuldig elk met $2^{3/T}$: thorium-232, $10.4 \times 2^{3/14.0} \approx
\qty{12.1}{TW}$; uranium-238, $9.6 \times 2^{3/4.468} \approx \qty{15.3}{TW}$;
kalium-40, $4.0 \times 2^{3/1.248} \approx \qty{21.2}{TW}$; uranium-235,
$0.42 \times 2^{3/0.704} \approx \qty{8.1}{TW}$. Totaal ongeveer \qty{56.7}{TW},
meer dan twee keer de \qty{24.4}{TW} van vandaag. De twee kortlevende isotopen leveren
$(21.2 + 8.1)/56.7 \approx \qty{52}{\percent}$ ervan, tegen
$4.42/24.4 \approx \qty{18}{\percent}$ vandaag.
\end{solution}

\begin{solution}{exo:g12:internal-heat:13}
Het vermogen is de massa per seconde maal de soortelijke warmte maal de
opwarming: $\num{3e8} = 1640 \times 4180 \times \Delta T$, dus $\Delta T =
\num{3e8}/(1640 \times 4180) \approx \qty{44}{K}$. Het gemiddelde continent verliest
\qty{0.0709}{W} per vierkante meter, dus is \qty{300}{MW} het natuurlijke warmteverlies
van $\num{3e8}/0.0709 \approx \qty{4.2e9}{m^2}$, ongeveer \qty{4200}{km^2}:
de centrale verzamelt warmte die de Aarde van nature over duizenden
vierkante kilometer laat ontsnappen.
\end{solution}

\begin{solution}{exo:g12:internal-heat:14}
(1) 100 miljoen jaar is $\num{2.6e11}/\num{1e8} = 2600$ keer
korter dan 260 miljard jaar; 200 miljoen jaar 1300 keer korter.
(2) Convectie voert de warmte mee met het bewegende gesteente: de tijd om de
mantel te doorkruisen is zijn dikte gedeeld door de snelheid, $\qty{2891}{km}$ bij een paar
centimeter per jaar, ongeveer honderd miljoen jaar. Geleiding heeft zo'n
drager niet: haar tijd groeit met het kwadraat van de afstand, en over duizenden
kilometers wordt die enorm.
\end{solution}

\begin{solution}{exo:g12:internal-heat:15}
(1) $10/22.3 \approx \qty{45}{\percent}$; de rest is de afkoeling van de
mantel zelf, die zijn eigen opgeslagen warmte verliest. (2) De grens ligt op een straal van $6371 - 2891 =
\qty{3480}{km}$; haar oppervlakte is $4\pi (\num{3.48e6})^2 \approx
\qty{1.52e14}{m^2}$, en $\num{1e13}/\num{1.52e14} \approx
\qty{0.066}{W/m^2}$, ongeveer \qty{66}{mW/m^2}, zo'n \qty{72}{\percent} van het
gemiddelde aan het oppervlak. (3) De warmte steekt de grens over door geleiding naar
de onderste laag van de mantel; opgewarmd wordt dit gesteente lichter dan de
mantel erboven en kan het opstijgen, waarbij het de warmte omhoog meevoert: zo
beginnen pluimen vermoedelijk, diep beneden.
\end{solution}

\begin{solution}{pb:g12:internal-heat:1}
\textbf{1.} $4\pi (\num{6.371e6})^2 \approx \qty{5.101e14}{m^2}$, en
$2.073 + 3.028 = 5.101$: ja.

\textbf{2.} $3.028/5.101 \approx \qty{59}{\percent}$, ongeveer drie vijfde.

\textbf{3.} De oceaanbodem is jong: hij wordt heet gemaakt bij de ruggen en
koelt af terwijl hij zich verwijdert, en jonge \omterm{def:g11:lithosphere-asthenosphere:lithosphere}{lithosfeer} verliest veel warmte.

\textbf{4.} $490/\sqrt{4} = \qty{245}{mW/m^2}$.

\textbf{5.} $2.72 \times 0.026 \approx \qty{0.0707}{W/m^2} =
\qty{70.7}{mW/m^2}$: net onder het gemiddelde, er praktisch aan gelijk.

\textbf{6.} $0.0709 \times \num{2.073e14} \approx \qty{14.7}{TW}$.

\textbf{7.} $0.1054 \times \num{3.028e14} \approx \qty{31.9}{TW}$.

\textbf{8.} $14.7 + 31.9 \approx \qty{46.6}{TW}$, de \qty{47}{TW} van de
gepubliceerde schatting (het wereldgemiddelde uit het hoofdstuk geeft \qty{46.7}{TW}, de
waarde die hieronder wordt gebruikt).

\textbf{9.} $\num{4.66e13}/\num{5.101e14} \approx \qty{0.0914}{W/m^2}$,
ongeveer \qty{91}{mW/m^2}.

\textbf{10.} $0.0243 \times \num{2.073e14} \approx \qty{5.0}{TW}$, een derde van
de \qty{14.7}{TW} die door de continenten verloren gaat: de radioactiviteit van de
\omterm{def:g12:continental-crust:continental-crust}{continentale korst} is een groot deel van hun \omterm{def:g12:internal-heat:heat-flow}{warmtestroom}.

\textbf{11.} Voor elke kilogram uranium bevat de Aarde \qty{3.9}{kg}
thorium, dat $3.9 \times 26.3 \approx \qty{102.6}{\micro W}$ vrijmaakt,
tegen \qty{94.4}{\micro W} voor het uranium-238. Uranium-238 maakt dus
$94.4/(94.4 + 102.6) \approx 0.479$ van de \qty{20.0}{TW}: ongeveer
\qty{9.6}{TW}. Uranium-235: $9.59 \times 4.10/94.4 \approx \qty{0.42}{TW}$.

\textbf{12.} $20.0 + 0.42 + 4 \approx \qty{24.4}{TW}$.

\textbf{13.} Graniet en de andere gesteenten van de \omterm{def:g12:continental-crust:continental-crust}{continentale korst}, ongeveer
twintig keer rijker dan de ultramafische gesteenten van de mantel.

\textbf{14.} Groter: elke isotoop was toen overvloediger, twee keer zoveel voor
elke halveringstijd terug; 3 miljard jaar geleden was het totaal meer dan twee keer
dat van vandaag.

\textbf{15.} $46.7 - 24.4 = \qty{22.3}{TW}$.

\textbf{16.} Uit de warmte die sinds het ontstaan van de Aarde is opgeslagen: de \omterm{def:g12:internal-heat:secular-cooling}{seculaire
afkoeling} van de mantel en de kern.

\textbf{17.} $22.3/46.7 \approx \qty{48}{\percent}$.

\textbf{18.} Van $11/46.7 \approx \qty{24}{\percent}$ tot $38/46.7 \approx
\qty{81}{\percent}$: de geoneutrino's vernauwen dit bereik sterk.

\textbf{19.} \omterm{def:g12:internal-heat:convection}{Mantelconvectie}, die het grootste deel van de warmte meevoert; alleen de laatste
tientallen kilometers worden door geleiding doorkruist.

\textbf{20.} $24.4/46.7 \approx 0.52$: radioactiviteit levert vandaag
\textbf{ongeveer de helft} van het warmteverlies van de Aarde, en het trage afkoelen van
de planeet de andere helft.
\end{solution}
